% GENERATED FILE. Edit canonical lessons, then run npm run generate:books.
% canonical-source-hash: fnv1a64:1fd426c59ade463d
% canonical-lessons: TE-C327-okappudu, TE-C327-raju, TE-C327-rani, TE-C327-rajyam, TE-C327-elu

\chapter{Once Upon a Time}
\label{ch:te-once-upon-a-time}
\hlchaptermodality{327}

\begin{chapteropening}
\textbf{By the end of this chapter:} \emph{I can say at one time, a king, a queen, a kingdom and to reign over.}

Complete the last lesson of chapter 327: I can say at one time, a king, a queen, a kingdom and to reign over.
\end{chapteropening}

\section[\texorpdfstring{\te{ఒకప్పుడు} (okappuḍu)}{ఒకప్పుడు (okappuḍu)}]{\te{ఒకప్పుడు} (okappuḍu) — at one time; once}

\label{lesson:TE-C327-okappudu}



\begin{quote}
Before the new one: say the Telugu for to steal, then the Telugu for to run away.
\end{quote}

\subsection*{You'll want to know: \te{ఒకప్పుడు}}
\textbf{\te{ఒకప్పుడు}} — \emph{okappuḍu} — \textquotedblleft{}at one time; once\textquotedblright{}.

\begin{quote}
\te{ఒకప్పుడు} \te{ఇక్కడ} \te{ఒక} \te{పెద్ద} \te{అడవి} \te{ఉండేది}.
\end{quote}

\begin{quote}
\emph{okappuḍu ikkaḍa oka pedda aḍavi uṇḍēdi.}
\end{quote}

\begin{quote}
At one time there was a big forest here.
\end{quote}

Stories often open with \te{ఒకప్పుడు}, the way English opens with once upon a time: \te{ఒక}, one, and \te{అప్పుడు}, then.

\begin{cousinweb}
\textbf{\te{ఒకప్పుడు}} (\emph{okappuḍu}) is a native Telugu compound: \emph{oka}, \textquotedblleft{}one\textquotedblright{} (DEDR 990), and \emph{appuḍu}, \textquotedblleft{}at that time\textquotedblright{} (DEDR 4559).
\end{cousinweb}

\begin{grammarlens}[title={What you've built}]
One more word for telling what happened.
\end{grammarlens}

\subsection*{Your turn}
\begin{itemize}
\raggedright
  \item \emph{Say it:} \emph{okappuḍu}
  \item \emph{Say it:} \emph{okappuḍu}, once more
  \item \emph{Say it:} \emph{okappuḍu ikkaḍa oka pedda aḍavi uṇḍēdi.} — At one time there was a big forest here.
  \item \emph{Recall:} say the Telugu for to steal, then the Telugu for to run away, then say \emph{okappuḍu} again
\end{itemize}

\begin{tcolorbox}[breakable,colback=teal!4,colframe=teal!35!black,title={Before you move on}]
What does \te{ఒకప్పుడు} mean? (\textquotedblleft{}At one time; once\textquotedblright{}.) Say it once more.
\end{tcolorbox}
\section[\texorpdfstring{\te{రాజు} (rāju)}{రాజు (rāju)}]{\te{రాజు} (rāju) — a king}

\label{lesson:TE-C327-raju}



\begin{quote}
Before the new one: say the Telugu for to run away, then the Telugu for at one time.
\end{quote}

\subsection*{You'll want to know: \te{రాజు}}
\textbf{\te{రాజు}} — \emph{rāju} — \textquotedblleft{}a king\textquotedblright{}.

\begin{quote}
\te{ఆ} \te{రాజు} \te{చాలా} \te{మంచివాడు}.
\end{quote}

\begin{quote}
\emph{ā rāju cālā mañcivāḍu.}
\end{quote}

\begin{quote}
That king is a very good man.
\end{quote}

Plural: \te{రాజులు}, \emph{rājulu}.

The capital, \te{రాజధాని}, is the seat of the king: the same Sanskrit word sits inside it.

\begin{cousinweb}
\textbf{\te{రాజు}} (\emph{rāju}) is borrowed from Sanskrit \emph{rājan}, \textquotedblleft{}a chieftain, a king\textquotedblright{} (CDIAL 10679). Telugu gives it the native ending -\te{ఉ}.
\end{cousinweb}

\begin{grammarlens}[title={What you've built}]
One more word for telling what happened.
\end{grammarlens}

\subsection*{Your turn}
\begin{itemize}
\raggedright
  \item \emph{Say it:} \emph{rāju}
  \item \emph{Say it:} \emph{rāju}, once more
  \item \emph{Say it:} \emph{ā rāju cālā mañcivāḍu.} — That king is a very good man.
  \item \emph{Recall:} say the Telugu for to run away, then the Telugu for at one time, then say \emph{rāju} again
\end{itemize}

\begin{tcolorbox}[breakable,colback=teal!4,colframe=teal!35!black,title={Before you move on}]
What does \te{రాజు} mean? (\textquotedblleft{}A king\textquotedblright{}.) Say it once more.
\end{tcolorbox}
\section[\texorpdfstring{\te{రాణి} (rāṇi)}{రాణి (rāṇi)}]{\te{రాణి} (rāṇi) — a queen}

\label{lesson:TE-C327-rani}



\begin{quote}
Before the new one: say the Telugu for at one time, then the Telugu for a king.
\end{quote}

\subsection*{You'll want to know: \te{రాణి}}
\textbf{\te{రాణి}} — \emph{rāṇi} — \textquotedblleft{}a queen\textquotedblright{}.

\begin{quote}
\te{రాణి} \te{రాజుతో} \te{అడవికి} \te{వెళ్ళింది}.
\end{quote}

\begin{quote}
\emph{rāṇi rājutō aḍaviki veḷḷindi.}
\end{quote}

\begin{quote}
The queen went to the forest with the king.
\end{quote}

\begin{cousinweb}
\textbf{\te{రాణి}} (\emph{rāṇi}) is borrowed from Sanskrit \emph{rājñī}, \textquotedblleft{}a queen\textquotedblright{} (CDIAL 10692), the feminine of \emph{rājan}, \textquotedblleft{}a king\textquotedblright{}.
\end{cousinweb}

\begin{grammarlens}[title={What you've built}]
One more word for telling what happened.
\end{grammarlens}

\subsection*{Your turn}
\begin{itemize}
\raggedright
  \item \emph{Say it:} \emph{rāṇi}
  \item \emph{Say it:} \emph{rāṇi}, once more
  \item \emph{Say it:} \emph{rāṇi rājutō aḍaviki veḷḷindi.} — The queen went to the forest with the king.
  \item \emph{Recall:} say the Telugu for at one time, then the Telugu for a king, then say \emph{rāṇi} again
\end{itemize}

\begin{tcolorbox}[breakable,colback=teal!4,colframe=teal!35!black,title={Before you move on}]
What does \te{రాణి} mean? (\textquotedblleft{}A queen\textquotedblright{}.) Say it once more.
\end{tcolorbox}
\section[\texorpdfstring{\te{రాజ్యం} (rājyaṁ)}{రాజ్యం (rājyaṁ)}]{\te{రాజ్యం} (rājyaṁ) — a kingdom}

\label{lesson:TE-C327-rajyam}



\begin{quote}
Before the new one: say the Telugu for a king, then the Telugu for a queen.
\end{quote}

\subsection*{You'll want to know: \te{రాజ్యం}}
\textbf{\te{రాజ్యం}} — \emph{rājyaṁ} — \textquotedblleft{}a kingdom\textquotedblright{}.

\begin{quote}
\te{ఆ} \te{రాజ్యంలో} \te{ప్రజలు} \te{సంతోషంగా} \te{ఉండేవారు}.
\end{quote}

\begin{quote}
\emph{ā rājyaṁlō prajalu santōṣaṅgā uṇḍēvāru.}
\end{quote}

\begin{quote}
In that kingdom the people lived happily.
\end{quote}

Plural: \te{రాజ్యాలు}, \emph{rājyālu}.

A Sanskrit -a noun ending in -\te{ం}: its plural is \te{రాజ్యాలు}, the way \te{పుస్తకం} gives \te{పుస్తకాలు}. In the old stories a kingdom depends on its king: \te{రాజ్యం} \te{రాజు} \te{మీద} \te{ఆధారపడి} \te{ఉంటుంది}.

\begin{cousinweb}
\textbf{\te{రాజ్యం}} (\emph{rājyaṁ}) is borrowed from Sanskrit \emph{rājya}, \textquotedblleft{}kingship\textquotedblright{}, and later \textquotedblleft{}a kingdom\textquotedblright{} (CDIAL 10694), built on \emph{rājan}, \textquotedblleft{}a king\textquotedblright{}.
\end{cousinweb}

\begin{grammarlens}[title={What you've built}]
One more word for telling what happened.
\end{grammarlens}

\subsection*{Your turn}
\begin{itemize}
\raggedright
  \item \emph{Say it:} \emph{rājyaṁ}
  \item \emph{Say it:} \emph{rājyaṁ}, once more
  \item \emph{Say it:} \emph{ā rājyaṁlō prajalu santōṣaṅgā uṇḍēvāru.} — In that kingdom the people lived happily.
  \item \emph{Recall:} say the Telugu for a king, then the Telugu for a queen, then say \emph{rājyaṁ} again
\end{itemize}

\begin{tcolorbox}[breakable,colback=teal!4,colframe=teal!35!black,title={Before you move on}]
What does \te{రాజ్యం} mean? (\textquotedblleft{}A kingdom\textquotedblright{}.) Say it once more.
\end{tcolorbox}
\section[\texorpdfstring{\te{ఏలు} (ēlu)}{ఏలు (ēlu)}]{\te{ఏలు} (ēlu) — to reign over}

\label{lesson:TE-C327-elu}



\begin{quote}
Before the new one: say the Telugu for a queen, then the Telugu for a kingdom.
\end{quote}

\subsection*{You'll want to know: \te{ఏలు}}
\textbf{\te{ఏలు}} — \emph{ēlu} — \textquotedblleft{}to reign over\textquotedblright{}.

\begin{quote}
\te{ఆ} \te{రాజు} \te{చాలా} \te{సంవత్సరాలు} \te{రాజ్యం} \te{ఏలాడు}.
\end{quote}

\begin{quote}
\emph{ā rāju cālā saṁvatsarālu rājyaṁ ēlāḍu.}
\end{quote}

\begin{quote}
That king ruled the kingdom for many years.
\end{quote}

Past (I …): \te{ఏలాను}, \emph{ēlānu}.

A native verb of the old stories: \te{రాజ్యం} \te{ఏలు}, to rule a kingdom.

\begin{cousinweb}
\textbf{\te{ఏలు}} (\emph{ēlu}) is native Dravidian: DEDR 5157 lists Telugu \emph{ēlu}, \textquotedblleft{}to rule, govern\textquotedblright{}, with Tamil \emph{āḷ}, \textquotedblleft{}to rule, reign over\textquotedblright{}.
\end{cousinweb}

\begin{grammarlens}[title={What you've built}]
One more word for telling what happened.
\end{grammarlens}

\subsection*{Your turn}
\begin{itemize}
\raggedright
  \item \emph{Say it:} \emph{ēlu}
  \item \emph{Say it:} \emph{ēlu}, once more
  \item \emph{Say it:} \emph{ā rāju cālā saṁvatsarālu rājyaṁ ēlāḍu.} — That king ruled the kingdom for many years.
  \item \emph{Recall:} say the Telugu for a queen, then the Telugu for a kingdom, then say \emph{ēlu} again
\end{itemize}

\begin{tcolorbox}[breakable,colback=teal!4,colframe=teal!35!black,title={Before you move on}]
What does \te{ఏలు} mean? (\textquotedblleft{}To reign over\textquotedblright{}.) Say it once more.
\end{tcolorbox}
